\section{Introduction} 
In software development, tests are essential for identifying defects and verifying that fixes work correctly \cite{kapfhammer2004software}. Ideally, test results remain stable across runs of an unchanged codebase. However, flaky tests, that is, tests that pass or fail nondeterministically without any change to the code or the execution environment, are a common reality \cite{pinto2020vocabulary}. Their causes are diverse, including concurrency issues, asynchronous operations, order dependencies between tests, improper management of external resources, dependence on the system clock, and the use of random values \cite{luo2014empirical}, as well as overly narrow ranges of expected values and timeout-related issues \cite{eck2019understanding}.
Although flaky tests do not directly affect software functionality, they impose a substantial cost on delivery pipelines. Developers report wasted effort analyzing failures, ignored failures, and loss of trust in the test suite \cite{eck2019understanding, gruber2022survey}, and at Google 16\% of tests exhibited flaky behavior, consuming up to 16\% of testing capacity in reruns \cite{micco2017state}. Detecting flakiness through reruns is expensive, since tests that fail rarely require many executions \cite{lam2020study}. Static prediction from the test code is therefore attractive: it requires running nothing and produces a human-readable explanation.
Pinto et al. \cite{pinto2020vocabulary} showed that flaky tests share a characteristic vocabulary, that is, identifiers and keywords that occur disproportionately in them, and achieved an F-score of 0.95 in Java. The idea was replicated for JavaScript \cite{soratto2023vocabulary}, extended to five languages \cite{ahmad2025understanding}, and applied to Swift \cite{medeiros2026vocabulary}. To the best of our knowledge, it has never been evaluated for C\#, even though the language has relevant particularities. Its concurrency model centers on the Task-based asynchronous pattern (Task, async/await, CancellationToken) and the thread pool, which introduces its own timing and scheduling dependencies. In tests from Microsoft .NET projects, concurrency is the leading cause of flakiness, with fixed waits and timeouts as a frequent pattern \cite{leesatapornwongsa2022flakerepro}. Test suites are written with several testing frameworks (xUnit, NUnit, MSTest) on two execution platforms, and libraries are often compiled for multiple target frameworks, so the same test may run under different runtimes. Integration suites also depend on containerized services and on environment behaviors such as locale, time zone, and file system paths. If the lexical patterns that signal flakiness in Java, JavaScript, and Swift do not transfer to C\#, vocabulary-based tools built around them would offer little help to .NET developers. This is the gap this work addresses.
To address it, we investigate static prediction of flaky tests in C\# from the vocabulary of test source code. Our approach, organized in six steps (Section 4), selects open-source C\# projects and systematically reruns their test suites, 100 times in an automated fashion and again 20 times under controlled isolation; labels each test as flaky, broken, or stable according to its outcomes; resolves inherited tests to the method that declares them; extracts the vocabulary and structural features of each test with the Roslyn compiler platform; and trains classifiers on this vocabulary to predict a test's propensity to be flaky without executing it, while also extracting an interpretable ranking of the tokens most strongly associated with flakiness.
We applied this methodology to 17 projects, with 19,211 test cases and 12,665 labeled test methods, of which 30 are flaky and 12,635 are stable (0.24\%). Among the eight classifiers evaluated, tree ensembles performed best under cross-validation, with Extra Trees reaching an MCC of 0.53 and an AUC of 0.86, and Random Forest an MCC of 0.51 and an AUC of 0.87. These results are well above random guessing (MCC 0.02), but only moderately above a structure-only baseline (MCC 0.36) and a vocabulary threshold rule (MCC 0.40). When each repository is held out and evaluated at its natural base rate (0.39\%), AUC drops to 0.75--0.78 and precision stays at 2--3\%; thus, the approach is better suited to ranking tests and prioritizing reruns than to gating a pipeline: the top quartile of each repository's ranking contains 23 of the 30 flaky tests.
The main contributions of this work are:
A dataset of flaky, broken, and stable C\# tests, built from 100 automated and 20 controlled reruns of 17 open-source projects, with inherited tests resolved to the methods that declare them.
An evaluation of eight classifiers for vocabulary-based static prediction in C\#, under repeated cross-validation and under cross-repository validation at the natural base rate, showing that the approach transfers to C\#, with clear limits.
A C\# flakiness vocabulary, which separates frequent general markers from rare, specific ones and shows that no robust marker of stability emerges.
An analysis of the model's errors in the test code, with sensitivity analyses without attributes and with helper-method expansion, which outline the limits of lexical prediction.
An assessment of label robustness, comparing the two rerun rounds (which agree on only 13 of the 36 flaky tests), and documentation of the pitfalls of rerun studies in .NET that distorted our initial counts.